\documentclass[12pt,a4paper]{report}
\usepackage[utf8]{inputenc}
\usepackage[italian]{babel}
\usepackage{geometry}
\usepackage{graphicx}
\usepackage{tikz}
\usepackage{listings}
\usepackage{xcolor}
\usepackage{hyperref}
\usepackage{booktabs}
\usepackage{amsmath}
\usepackage{amssymb}
\usepackage{fancyhdr}
\usepackage{titlesec}
\usepackage{longtable}
\usepackage{array}
\usepackage{caption}
\usepackage{subcaption}
\usepackage{float}
\usepackage{tcolorbox}
\usepackage{enumitem}

\usetikzlibrary{shapes.geometric, arrows.meta, positioning, fit, backgrounds, calc, decorations.pathreplacing}

\geometry{margin=2.5cm}

\definecolor{codebg}{RGB}{248,248,248}
\definecolor{codeframe}{RGB}{200,200,200}
\definecolor{darkblue}{RGB}{0,0,139}
\definecolor{darkgreen}{RGB}{0,100,0}
\definecolor{darkred}{RGB}{139,0,0}
\definecolor{accent}{RGB}{70,130,180}

\lstset{
    backgroundcolor=\color{codebg},
    frame=single,
    rulecolor=\color{codeframe},
    basicstyle=\ttfamily\small,
    keywordstyle=\color{darkblue}\bfseries,
    commentstyle=\color{darkgreen}\itshape,
    stringstyle=\color{darkred},
    numbers=left,
    numberstyle=\tiny\color{gray},
    numbersep=8pt,
    breaklines=true,
    breakatwhitespace=true,
    showstringspaces=false,
    tabsize=2,
    captionpos=b,
    language=JavaScript,
    morekeywords={async,await,export,import,from,class,extends,constructor}
}

\hypersetup{
    colorlinks=true,
    linkcolor=darkblue,
    urlcolor=accent,
    citecolor=darkgreen,
    bookmarks=true,
    bookmarksopen=true
}

\pagestyle{fancy}
\fancyhf{}
\fancyhead[L]{\leftmark}
\fancyhead[R]{\thepage}
\renewcommand{\headrulewidth}{0.4pt}

\titleformat{\chapter}[display]
    {\normalfont\huge\bfseries}
    {\chaptertitlename\ \thechapter}{20pt}{\Huge}
\titleformat{\section}
    {\normalfont\Large\bfseries}
    {\thesection}{1em}{}
\titleformat{\subsection}
    {\normalfont\large\bfseries}
    {\thesubsection}{1em}{}

\newtcolorbox{infobox}[1][]{
    colback=blue!5!white,
    colframe=accent,
    fonttitle=\bfseries,
    title=#1
}

\newtcolorbox{warnbox}[1][]{
    colback=orange!5!white,
    colframe=orange!75!black,
    fonttitle=\bfseries,
    title=#1
}

\newtcolorbox{secbox}[1][]{
    colback=green!5!white,
    colframe=darkgreen,
    fonttitle=\bfseries,
    title=#1
}

\begin{document}

%==============================================================================
% TITOLO
%==============================================================================
\begin{titlepage}
    \centering
    \vspace*{2cm}
    {\Huge\bfseries Sistema di Autenticazione\\[0.5cm] ExtraTracker}\\[2cm]
    {\LARGE Documentazione Tecnica Completa}\\[1cm]
    {\large Architettura, Sicurezza e Implementazione}\\[3cm]
    
    \begin{tikzpicture}
        \node[draw=accent, line width=2pt, rounded corners=10pt, inner sep=20pt] {
            \begin{minipage}{0.7\textwidth}
                \centering
                \textbf{Stack Tecnologico}\\[0.5cm]
                Node.js $\bullet$ Express $\bullet$ MongoDB $\bullet$ React\\
                Argon2id $\bullet$ JWT $\bullet$ Redis $\bullet$ 2FA/TOTP
            \end{minipage}
        };
    \end{tikzpicture}
    
    \vfill
    {\large Versione 2.0 --- \today}
\end{titlepage}

%==============================================================================
% INDICE
%==============================================================================
\tableofcontents
\newpage

%==============================================================================
% CAPITOLO 1: OVERVIEW ARCHITETTURALE
%==============================================================================
\chapter{Overview Architetturale}

\section{Visione Generale del Sistema}

Il sistema di autenticazione di ExtraTracker rappresenta un'implementazione enterprise-grade che integra meccanismi di sicurezza multi-layer, audit completo delle operazioni e gestione avanzata delle sessioni. L'architettura è progettata secondo i principi\footnote{OWASP Authentication Cheat Sheet, NIST Digital Identity Guidelines}:

\begin{itemize}[leftmargin=*]
    \item \textbf{Defense in Depth}: Multiple layer di protezione sovrapposti
    \item \textbf{Principle of Least Privilege}: Accesso minimo necessario
    \item \textbf{Fail Secure}: Gestione sicura degli errori
    \item \textbf{Auditability}: Tracciabilità completa di ogni azione
\end{itemize}

\section{Componenti Principali}

\begin{figure}[H]
    \centering
    \begin{tikzpicture}[
        node distance=1.5cm and 2cm,
        box/.style={rectangle, draw=accent, thick, rounded corners, minimum width=3cm, minimum height=1cm, align=center, fill=accent!10},
        db/.style={cylinder, shape border rotate=90, draw=darkgreen, thick, aspect=0.25, minimum width=2.5cm, minimum height=1cm, align=center, fill=darkgreen!10},
        arrow/.style={-{Stealth[length=3mm]}, thick}
    ]
        % Client
        \node[box, fill=blue!20] (client) {\textbf{Client React}\\Browser/SPA};
        
        % Middleware
        \node[box, right=of client] (proxy) {\textbf{Vite DevServer}\\Proxy API};
        
        % Backend layers
        \node[box, right=of proxy, fill=orange!20] (security) {\textbf{Security Layer}\\Helmet, CORS, RateLimit};
        
        \node[box, below=of security, fill=red!20] (csrf) {\textbf{CSRF Protection}\\Double-submit Cookie};
        
        \node[box, below=of csrf, fill=purple!20] (jwt) {\textbf{JWT Validation}\\Token Verification};
        
        % Controllers
        \node[box, right=of security, fill=yellow!20] (auth) {\textbf{Auth Controller}\\Login/Logout/Refresh};
        
        % Services
        \node[box, below=of auth, fill=cyan!20] (service) {\textbf{Auth Service}\\Business Logic};
        
        % Audit
        \node[box, below=of service, fill=gray!20] (audit) {\textbf{Audit Service}\\Security Logging};
        
        % Database
        \node[db, right=of auth] (mongo) {\textbf{MongoDB}\\User Collection};
        \node[db, below=of mongo] (auditdb) {\textbf{MongoDB}\\Audit Collection};
        \node[db, below=of auditdb] (redis) {\textbf{Redis}\\Rate Limit/Blacklist};
        
        % Arrows
        \draw[arrow] (client) -- (proxy);
        \draw[arrow] (proxy) -- (security);
        \draw[arrow] (security) -- (csrf);
        \draw[arrow] (csrf) -- (jwt);
        \draw[arrow] (jwt) -- (auth);
        \draw[arrow] (auth) -- (service);
        \draw[arrow] (service) -- (audit);
        \draw[arrow] (service) -- (mongo);
        \draw[arrow] (audit) -- (auditdb);
        \draw[arrow] (security) -- ++(0,1.5) -| (redis);
        
    \end{tikzpicture}
    \caption{Architettura del Sistema di Autenticazione}
    \label{fig:architecture}
\end{figure}

\section{Stack Tecnologico}

\subsection{Backend (Node.js/Express)}

\begin{longtable}{@{}p{4cm}p{5cm}p{5cm}@{}}
\toprule
\textbf{Componente} & \textbf{Libreria} & \textbf{Scopo} \\
\midrule
\endhead
Password Hashing & argon2 & Argon2id (winner PHC) \\
JWT Handling & jsonwebtoken (jose) & Token signing/verification \\
2FA/TOTP & otplib & RFC 6238 compliant \\
QR Generation & qrcode & Setup 2FA \\
Device Parsing & ua-parser-js & Fingerprinting \\
GeoIP & geoip-lite & Location tracking \\
Encryption & crypto (native) & AES-256-GCM \\
Rate Limiting & express-rate-limit & DoS protection \\
Security Headers & helmet & OWASP headers \\
CSRF Protection & csurf & Double-submit pattern \\
Input Validation & express-validator & Schema validation \\
\bottomrule
\end{longtable}

\subsection{Database Layer}

\begin{infobox}[MongoDB Collections]
\begin{itemize}
    \item \texttt{users}: Documenti utente con embedded arrays
    \item \texttt{auditlogs}: Log di sicurezza con TTL index
    \item Redis (opzionale): Rate limiting, session blacklist
\end{itemize}
\end{infobox}

%==============================================================================
% CAPITOLO 2: FLUSSO DI AUTENTICAZIONE
%==============================================================================
\chapter{Flusso di Autenticazione}

\section{Sequence Diagram Completo}

\begin{figure}[H]
    \centering
    \begin{tikzpicture}[
        scale=0.8,
        transform shape
    ]
        % Partecipanti
        \node[draw, fill=blue!20, minimum width=2.5cm] (client) at (0,0) {\textbf{Client}};
        \node[draw, fill=orange!20, minimum width=2.5cm] (proxy) at (4,0) {\textbf{Proxy}};
        \node[draw, fill=green!20, minimum width=2.5cm] (api) at (8,0) {\textbf{API Server}};
        \node[draw, fill=purple!20, minimum width=2.5cm] (auth) at (12,0) {\textbf{AuthService}};
        \node[draw, fill=red!20, minimum width=2.5cm] (db) at (16,0) {\textbf{Database}};
        
        % Lifelines
        \draw[dashed] (client) -- (0,-22);
        \draw[dashed] (proxy) -- (4,-22);
        \draw[dashed] (api) -- (8,-22);
        \draw[dashed] (auth) -- (12,-22);
        \draw[dashed] (db) -- (16,-22);
        
        % Messaggi
        \draw[-{Stealth}] (0,-1) -- (4,-1) node[midway, above, sloped] {\small 1. GET /csrf};
        \draw[-{Stealth}] (4,-1.5) -- (8,-1.5) node[midway, above, sloped] {\small 2. forward};
        \draw[-{Stealth}] (8,-2) -- (4,-2) node[midway, above, sloped] {\small 3. Set-Cookie};
        \draw[-{Stealth}] (4,-2.5) -- (0,-2.5) node[midway, above, sloped] {\small 4. csrfToken};
        
        \draw[-{Stealth}, thick] (0,-4) -- (4,-4) node[midway, above, sloped] {\small 5. POST /login + CSRF};
        \draw[-{Stealth}] (4,-4.5) -- (8,-4.5) node[midway, above, sloped] {\small 6. forward};
        
        \draw[-{Stealth}] (8,-5.5) -- (12,-5.5) node[midway, above, sloped] {\small 7. login(credentials)};
        \draw[-{Stealth}] (12,-6) -- (16,-6) node[midway, above, sloped] {\small 8. findForLogin()};
        \draw[-{Stealth}] (16,-6.5) -- (12,-6.5) node[midway, above, sloped] {\small 9. User document};
        
        \draw[-{Stealth}] (12,-7.5) -- (12,-7.8) node[right] {\small 10. argon2.verify()};
        
        \draw[-{Stealth}] (12,-8.5) -- (16,-8.5) node[midway, above, sloped] {\small 11. Audit log};
        
        \draw[-{Stealth}] (12,-9.5) -- (12,-10.5) node[right, align=left] {\small 12. generateAccessToken()\\\small \quad iss, aud, sub, jti, dfp};
        
        \draw[-{Stealth}] (12,-11) -- (12,-12) node[right, align=left] {\small 13. generateRefreshToken()\\\small \quad hash, jti, encrypt};
        
        \draw[-{Stealth}] (12,-12.5) -- (16,-12.5) node[midway, above, sloped] {\small 14. Update session array};
        
        \draw[-{Stealth}] (12,-13.5) -- (8,-13.5) node[midway, above, sloped] {\small 15. Tokens + User};
        \draw[-{Stealth}] (8,-14) -- (4,-14) node[midway, above, sloped] {\small 16. Set-Cookie (2)};
        \draw[-{Stealth}, thick] (4,-14.5) -- (0,-14.5) node[midway, above, sloped] {\small 17. 200 OK + user data};
        
        % Subsequent request
        \draw[-{Stealth}, thick, dashed] (0,-16) -- (4,-16) node[midway, above, sloped] {\small 18. GET /api/resource + Cookie};
        \draw[-{Stealth}, dashed] (4,-16.5) -- (8,-16.5) node[midway, above, sloped] {\small 19. forward};
        \draw[-{Stealth}, dashed] (8,-17) -- (12,-17) node[midway, above, sloped] {\small 20. verifyToken()};
        \draw[-{Stealth}, dashed] (12,-17.5) -- (8,-17.5) node[midway, above, sloped] {\small 21. Payload valid};
        \draw[-{Stealth}, dashed] (8,-18) -- (4,-18) node[midway, above, sloped] {\small 22. 200 + data};
        \draw[-{Stealth}, thick, dashed] (4,-18.5) -- (0,-18.5) node[midway, above, sloped] {\small 23. Response};
        
        % Activation boxes
        \fill[blue!10] (-0.3,-1) rectangle (0.3,-2.5);
        \fill[blue!10] (-0.3,-4) rectangle (0.3,-14.5);
        \fill[blue!10] (-0.3,-16) rectangle (0.3,-18.5);
        
        \fill[green!10] (7.7,-1.5) rectangle (8.3,-2);
        \fill[green!10] (7.7,-4.5) rectangle (8.3,-14);
        \fill[green!10] (7.7,-16.5) rectangle (8.3,-18);
        
        \fill[purple!10] (11.7,-5.5) rectangle (12.3,-13.5);
        \fill[purple!10] (11.7,-17) rectangle (12.3,-17.5);
        
        \fill[red!10] (15.7,-6) rectangle (16.3,-6.5);
        \fill[red!10] (15.7,-8.5) rectangle (16.3,-8.8);
        \fill[red!10] (15.7,-12.5) rectangle (16.3,-12.8);
        
    \end{tikzpicture}
    \caption{Sequence Diagram: Login Flow Completo}
    \label{fig:sequence}
\end{figure}

\section{Fasi Dettagliate}

\subsection{Fase 1: CSRF Token Acquisition}

Prima di qualsiasi operazione di mutazione, il client deve ottenere un token CSRF:

\begin{lstlisting}[caption={Endpoint CSRF}, label={lst:csrf}]
GET /api/auth/csrf
Response: { success: true, data: { csrfToken: "..." } }
Set-Cookie: csrfToken=<random>; HttpOnly=false; Secure; SameSite=Strict
\end{lstlisting}

Il token CSRF implementa il pattern\footnote{OWASP Cross-Site Request Forgery Prevention}:
\begin{itemize}
    \item \textbf{Double-Submit Cookie}: Token diverso dal session cookie
    \item \textbf{Synchronizer Token}: Server-side validation
    \item \textbf{SameSite=Strict}: Protezione contro CSRF cross-origin
\end{itemize}

\subsection{Fase 2: Credential Submission}

\begin{lstlisting}[caption={Login Request}, label={lst:login}]
POST /api/auth/login
Headers:
  Content-Type: application/json
  X-CSRF-Token: <csrf_token>
  Cookie: csrfToken=<csrf_cookie>

Body: {
  email: "user@example.com",
  password: "user_password",
  twoFactorCode?: "123456"
}
\end{lstlisting}

\subsection{Fase 3: Server-Side Processing}

\begin{figure}[H]
    \centering
    \begin{tikzpicture}[
        node distance=0.8cm,
        process/.style={rectangle, draw, fill=blue!20, minimum width=10cm, minimum height=0.8cm, align=center},
        decision/.style={diamond, draw, fill=yellow!20, aspect=2, align=center}
    ]
        \node[process] (p1) {1. Validazione input (express-validator)};
        \node[process, below=of p1] (p2) {2. Verifica CSRF token};
        \node[process, below=of p2] (p3) {3. Rate limiting check (Redis/memory)};
        \node[process, below=of p3] (p4) {4. User.findForLogin(email) - Proiezione campo password};
        \node[process, below=of p4] (p5) {5. Verifica account locked};
        \node[process, below=of p5] (p6) {6. argon2.verify(hash, password)};
        \node[process, below=of p6] (p7) {7. Gestione failed attempts};
        \node[process, below=of p7] (p8) {8. Verifica 2FA (se abilitato)};
        \node[process, below=of p8] (p9) {9. Generazione accessToken (JWT)};
        \node[process, below=of p9] (p10) {10. Generazione refreshToken + sessionData};
        \node[process, below=of p10] (p11) {11. Atomic update sessioni utente};
        \node[process, below=of p11] (p12) {12. Audit logging (async)};
        \node[process, below=of p12] (p13) {13. Suspicious activity detection};
        \node[process, below=of p13] (p14) {14. Set cookie HttpOnly Secure SameSite};
        
        \draw[-{Stealth}] (p1) -- (p2);
        \draw[-{Stealth}] (p2) -- (p3);
        \draw[-{Stealth}] (p3) -- (p4);
        \draw[-{Stealth}] (p4) -- (p5);
        \draw[-{Stealth}] (p5) -- (p6);
        \draw[-{Stealth}] (p6) -- (p7);
        \draw[-{Stealth}] (p7) -- (p8);
        \draw[-{Stealth}] (p8) -- (p9);
        \draw[-{Stealth}] (p9) -- (p10);
        \draw[-{Stealth}] (p10) -- (p11);
        \draw[-{Stealth}] (p11) -- (p12);
        \draw[-{Stealth}] (p12) -- (p13);
        \draw[-{Stealth}] (p13) -- (p14);
    \end{tikzpicture}
    \caption{Server-Side Login Processing Pipeline}
\end{figure}

%==============================================================================
% CAPITOLO 3: SICUREZZA PASSWORD
%==============================================================================
\chapter{Sicurezza delle Password}

\section{Argon2id: Password Hashing}

Il sistema utilizza \textbf{Argon2id}, vincitore della Password Hashing Competition (PHC) 2015.

\subsection{Configurazione Argon2}

\begin{lstlisting}[caption={Argon2 Configuration}, label={lst:argon2}]
const argon2Config = {
    type: argon2.argon2id,      // Ibrido: resistance a side-channel e GPU
    memoryCost: 65536,          // 64 MB di RAM
    timeCost: 3,                // 3 iterazioni
    parallelism: 4,             // 4 thread paralleli
    hashLength: 32              // 256 bit output
};
\end{lstlisting}

\begin{secbox}[Parametri Argon2id]
\begin{itemize}
    \item \textbf{Memory Cost (64MB)}: Mitigazione attacchi hardware custom (FPGA/ASIC)
    \item \textbf{Time Cost (3)}: Bilanciamento tra sicurezza e UX
    \item \textbf{Parallelism (4)}: Utilizzo multi-core moderni
    \item \textbf{Salt}: 16 bytes random generati automaticamente da argon2
\end{itemize}
\end{secbox}

\subsection{Formato Hash}

\begin{lstlisting}[caption={Esempio Hash Argon2id}]
$argon2id$v=19$m=65536,t=3,p=4$
    c29tZXNhbHR...base64...$        // Salt
    RGFJdfgzhj23...base64...        // Hash output
\end{lstlisting}

\section{Password History}

Per prevenire il riutilizzo delle password, il sistema mantiene uno storico criptato:

\begin{lstlisting}[caption={Schema Password History}, label={lst:passhistory}]
passwordHistory: [{
    hash: String,           // Hash precedente (argon2)
    changedAt: Date,        // Timestamp cambio
    reason: String          // CHANGE | RESET | ADMIN
}]
\end{lstlisting}

\begin{itemize}
    \item Massimo 5 password memorizzate (circular buffer)
    \item Confronto tramite \texttt{argon2.verify()} contro tutto lo storico
    \item Campo \texttt{reason} per audit trail completo
\end{itemize}

%==============================================================================
% CAPITOLO 4: JSON WEB TOKENS
%==============================================================================
\chapter{JSON Web Tokens (JWT)}

\section{Struttura Token}

\subsection{Access Token}

\begin{table}[H]
\centering
\begin{tabular}{@{}lll@{}}
\toprule
\textbf{Claim} & \textbf{Valore} & \textbf{Scopo} \\
\midrule
iss (Issuer) & ``silvi-api'' & Identifica il server \\
aud (Audience) & ``silvi-app'' & Destinatario valido \\
sub (Subject) & user\_id & Identificatore utente \\
jti (JWT ID) & UUID v4 & Unicità token \\
type & ``access'' & Distinguere da refresh \\
email & user@email & Denormalizzazione \\
dfp & SHA256(...) & Device fingerprint \\
iat & timestamp & Issued at \\
exp & timestamp & Expiration (15 min) \\
\bottomrule
\end{tabular}
\caption{Claims Access Token}
\end{table}

\subsection{Refresh Token}

\begin{lstlisting}[caption={Refresh Token Structure}]
{
    "iss": "silvi-api",
    "aud": "silvi-app",
    "sub": "user_id",
    "type": "refresh",
    "jti": "uuid",
    "iat": 1234567890,
    "exp": 1237170690  // 7 giorni
}
\end{lstlisting}

\section{Gestione Sessioni}

\begin{figure}[H]
    \centering
    \begin{tikzpicture}[
        box/.style={rectangle, draw, fill=blue!10, minimum width=4cm, minimum height=0.8cm, align=center},
        arrow/.style={-{Stealth}, thick}
    ]
        \node[box] (user) at (0,0) {\textbf{User Document}};
        \node[box] (session) at (6,2) {\textbf{Session Array}};
        \node[box] (grace) at (6,-2) {\textbf{Grace Period Array}};
        
        \node[rectangle, draw, fill=green!10, minimum width=3cm] (s1) at (12,3) {Session 1\\jti: abc\\device: Chrome};
        \node[rectangle, draw, fill=green!10, minimum width=3cm] (s2) at (12,1) {Session 2\\jti: def\\device: Mobile};
        
        \node[rectangle, draw, fill=orange!10, minimum width=3cm] (g1) at (12,-1) {Grace\\jti: old\\expired};
        \node[rectangle, draw, fill=orange!10, minimum width=3cm] (g2) at (12,-3) {Grace\\jti: prev\\rotation};
        
        \draw[arrow] (user) -- (session);
        \draw[arrow] (user) -- (grace);
        \draw[arrow] (session) -- (s1);
        \draw[arrow] (session) -- (s2);
        \draw[arrow] (grace) -- (g1);
        \draw[arrow] (grace) -- (g2);
        
    \end{tikzpicture}
    \caption{Struttura Sessioni Utente}
\end{figure}

\subsection{Refresh Token Rotation}

Ogni utilizzo di un refresh token genera:
\begin{enumerate}
    \item Nuovo refresh token con nuovo JTI
    \item Vecchio token spostato in \textit{grace period} (24h)
    \item Accesso con vecchio token =grace period= $\rightarrow$ nuovo token
    \item Accesso con token precedente al grace period $\rightarrow$ \textbf{revoca tutte le sessioni}
\end{enumerate}

\begin{warnbox}[Token Reuse Detection]
Se un refresh token già utilizzato (non in grace period) viene riutilizzato, il sistema considera questo comportamento come compromissione e revoca \textbf{tutte} le sessioni dell'utente.
\end{warnbox}

%==============================================================================
% CAPITOLO 5: TWO-FACTOR AUTHENTICATION
%==============================================================================
\chapter{Two-Factor Authentication (2FA)}

\section{TOTP: Time-based One-Time Password}

Implementazione conforme RFC 6238 con algoritmo HMAC-SHA1.

\subsection{Setup 2FA Flow}

\begin{figure}[H]
    \centering
    \begin{tikzpicture}[
        node distance=1.5cm,
        step/.style={rectangle, draw, fill=blue!10, minimum width=8cm, minimum height=1cm, align=center}
    ]
        \node[step] (s1) {1. POST /2fa/setup (requires auth)};
        \node[step, below=of s1] (s2) {2. Generazione secret (authenticator.generateSecret())};
        \node[step, below=of s2] (s3) {3. Generazione 10 backup codes};
        \node[step, below=of s3] (s4) {4. Hash backup codes (argon2)};
        \node[step, below=of s4] (s5) {5. Generazione QR code (keyuri + qrcode)};
        \node[step, below=of s5] (s6) {6. Response: secret, qrUrl, backupCodes (plaintext)};
        \node[step, below=of s6] (s7) {7. User salva backup codes in posto sicuro};
        \node[step, below=of s7] (s8) {8. POST /2fa/verify (code TOTP)};
        \node[step, below=of s8] (s9) {9. Abilitazione 2FA su account};
        
        \draw[-{Stealth}] (s1) -- (s2);
        \draw[-{Stealth}] (s2) -- (s3);
        \draw[-{Stealth}] (s3) -- (s4);
        \draw[-{Stealth}] (s4) -- (s5);
        \draw[-{Stealth}] (s5) -- (s6);
        \draw[-{Stealth}] (s6) -- (s7);
        \draw[-{Stealth}] (s7) -- (s8);
        \draw[-{Stealth}] (s8) -- (s9);
    \end{tikzpicture}
    \caption{2FA Setup Flow}
\end{figure}

\subsection{Formato URI (Key URI)}

\begin{lstlisting}[caption={OTP Auth URI}]
otpauth://totp/Silvi%20AI:email@example.com
    ?secret=JBSWY3DPEHPK3PXP
    &issuer=Silvi%20AI
    &algorithm=SHA1
    &digits=6
    &period=30
\end{lstlisting}

\subsection{Backup Codes}

\begin{itemize}
    \item 10 codici monouso, formato: \texttt{XXXX-XXXX}
    \item Hash Argon2 prima del salvataggio
    \item Utilizzo: confronto hash + rimozione immediata
    \item Regenerazione completa ad ogni setup 2FA
\end{itemize}

\section{Login con 2FA}

\begin{figure}[H]
    \centering
    \begin{tikzpicture}[
        node distance=1cm,
        process/.style={rectangle, draw, fill=blue!10, minimum width=6cm, minimum height=0.8cm},
        decision/.style={diamond, draw, fill=yellow!20, aspect=2, minimum width=3cm}
    ]
        \node[process] (p1) {Login con email/password};
        \node[decision, below=of p1] (d1) {2FA enabled?};
        \node[process, right=2cm of d1] (p2) {Return tokens};
        \node[process, below=of d1] (p3) {Return tempToken (5min)};
        \node[process, below=of p3] (p4) {Client richiede codice 2FA};
        \node[process, below=of p4] (p5) {POST /login con 2FA code};
        \node[decision, below=of p5] (d2) {Code valid?};
        \node[process, right=2cm of d2] (p6) {Audit fail};
        \node[process, below=of d2] (p7) {Return final tokens};
        
        \draw[-{Stealth}] (p1) -- (d1);
        \draw[-{Stealth}] (d1) -- node[above] {No} (p2);
        \draw[-{Stealth}] (d1) -- node[left] {Yes} (p3);
        \draw[-{Stealth}] (p3) -- (p4);
        \draw[-{Stealth}] (p4) -- (p5);
        \draw[-{Stealth}] (p5) -- (d2);
        \draw[-{Stealth}] (d2) -- node[above] {No} (p6);
        \draw[-{Stealth}] (d2) -- node[left] {Yes} (p7);
    \end{tikzpicture}
    \caption{2FA Login Flow}
\end{figure}

%==============================================================================
% CAPITOLO 6: DEVICE FINGERPRINTING
%==============================================================================
\chapter{Device Fingerprinting}

\section{Generazione Fingerprint}

Il fingerprint del dispositivo è calcolato combinando:

\begin{lstlisting}[caption={Device Fingerprint Generation}]
function generateDeviceFingerprint(req) {
    const components = [
        req.headers['user-agent'] || '',
        req.headers['accept-language'] || '',
        req.headers['accept-encoding'] || '',
        req.headers['dnt'] || '',
        req.ip
    ];
    
    const data = components.join('|');
    return crypto.createHash('sha256').update(data).digest('hex');
}
\end{lstlisting}

\section{Trusted Devices}

\begin{lstlisting}[caption={Trusted Device Schema}]
trustedDevices: [{
    fingerprint: String,      // SHA256 hash
    name: String,             // "Chrome on Windows"
    firstSeen: Date,
    lastUsed: Date,
    ip: String,               // Ultimo IP noto
    isTrusted: Boolean        // true = salta 2FA
}]
\end{lstlisting}

\subsection{Policy Trust Device}

\begin{itemize}
    \item Max 5 trusted devices per utente
    \item Revoca manuale possibile
    \item Auto-revoca dopo 90 giorni di inattività
    \item Device nuovo $\rightarrow$ richiede 2FA anche se ``Remember device''
\end{itemize}

%==============================================================================
% CAPITOLO 7: AUDIT LOGGING
%==============================================================================
\chapter{Audit Logging}

\section{Struttura Log di Sicurezza}

\begin{lstlisting}[caption={Audit Log Schema}]
{
    userId: ObjectId,
    userEmail: String,
    action: Enum[LOGIN_SUCCESS, LOGIN_FAILURE, ...],
    description: String,
    ip: String,
    userAgent: String,
    deviceFingerprint: String,
    location: {
        country: String,
        city: String,
        region: String,
        timezone: String,
        coordinates: { lat: Number, lon: Number }
    },
    success: Boolean,
    failureReason: String,
    sessionId: String,        // JTI del token
    severity: Enum[LOW, MEDIUM, HIGH, CRITICAL],
    metadata: Mixed,
    createdAt: Date           // TTL Index: 2 anni
}
\end{lstlisting}

\section{Eventi Tracciati}

\begin{longtable}{@{}lll@{}}
\toprule
\textbf{Evento} & \textbf{Severità} & \textbf{Trigger} \\
\midrule
\endhead
LOGIN\_SUCCESS & LOW & Login completato \\
LOGIN\_FAILURE & MEDIUM/HIGH & Password/errata \\
LOGOUT & LOW & Logout utente \\
LOGOUT\_ALL & HIGH & Revoca tutte sessioni \\
TOKEN\_REFRESH & LOW & Rotazione token \\
TOKEN\_REVOKED & HIGH & Rilevamento riuso \\
PASSWORD\_CHANGE & CRITICAL & Cambio password \\
PASSWORD\_RESET\_REQUEST & HIGH & Richiesta reset \\
PASSWORD\_RESET\_COMPLETE & CRITICAL & Reset completato \\
2FA\_ENABLED & CRITICAL & Abilitazione 2FA \\
2FA\_DISABLED & CRITICAL & Disabilitazione 2FA \\
2FA\_VERIFICATION & MEDIUM & Verifica codice \\
2FA\_BACKUP\_USED & HIGH & Uso backup code \\
REGISTER & MEDIUM & Nuovo account \\
ACCOUNT\_LOCKED & HIGH & Troppi tentativi \\
SUSPICIOUS\_ACTIVITY & CRITICAL & Pattern anomalo \\
DEVICE\_TRUSTED & MEDIUM & Nuovo device trusted \\
DEVICE\_REVOKED & HIGH & Revoca trust \\
\bottomrule
\end{longtable}

\section{Geolocalizzazione}

Utilizzo libreria \texttt{geoip-lite} per lookup IP:\\ https://github.com/geoip-lite/node-geoip

\begin{lstlisting}[caption={Geolocation Lookup}]
const geo = geoip.lookup(ip);
// Returns: { country, region, city, ll: [lat, lon], timezone }
\end{lstlisting}

\subsection{Privacy Considerazioni}

\begin{itemize}
    \item IP locali (127.0.0.1, ::1, 192.168.x.x) non tracciati
    \item TTL 2 anni per compliance GDPR
    \item Coordinate GPS arrotondate alla città (non indirizzo)
\end{itemize}

\section{Suspicious Activity Detection}

\begin{lstlisting}[caption={Detection Algorithm}]
async detectSuspiciousActivity(userId, currentIp, fingerprint) {
    const criteria = {
        multipleIps24h: count(unique IPs in 24h) > 3,
        multipleCountries: count(unique countries in 7d) > 2,
        newDevice: !fingerprint in trustedDevices
    };
    
    return {
        isSuspicious: criteria.multipleIps24h || 
                      criteria.multipleCountries,
        riskFactors: criteria,
        details: { uniqueIpsCount, uniqueCountriesCount }
    };
}
\end{lstlisting}

%==============================================================================
% CAPITOLO 8: RATE LIMITING
%==============================================================================
\chapter{Rate Limiting}

\section{Strategie di Rate Limiting}

\subsection{General Rate Limit}

\begin{lstlisting}[caption={General Limiter Configuration}]
{
    windowMs: 15 * 60 * 1000,    // 15 minuti
    max: 1000,                    // 1000 richieste
    keyGenerator: (req) => req.ip,
    handler: (req, res) => {
        throw AppError.tooManyRequests();
    }
}
\end{lstlisting}

\subsection{Auth-Specific Rate Limit}

\begin{lstlisting}[caption={Auth Limiter Configuration}]
{
    windowMs: 60 * 60 * 1000,     // 1 ora
    max: 5,                       // 5 tentativi login
    skipSuccessfulRequests: true,  // Non contare login OK
    keyGenerator: (req) => req.body.email || req.ip
}
\end{lstlisting}

\subsection{Strict Rate Limit}

Per endpoint sensibili (2FA setup, password reset):

\begin{lstlisting}[caption={Strict Limiter}]
{
    windowMs: 60 * 60 * 1000,
    max: 3,
    skipSuccessfulRequests: false
}
\end{lstlisting}

\section{Storage Backend}

\begin{itemize}
    \item \textbf{Redis} (produzione): Store distribuito, persistenza
    \item \textbf{Memory} (sviluppo): Map in-memory, reset su restart
\end{itemize}

%==============================================================================
% CAPITOLO 9: CSRF PROTECTION
%==============================================================================
\chapter{CSRF Protection}

\section{Double-Submit Cookie Pattern}

Il sistema implementa il pattern\footnote{OWASP CSRF Prevention} double-submit cookie:

\begin{figure}[H]
    \centering
    \begin{tikzpicture}[
        node distance=2cm,
        entity/.style={rectangle, draw, minimum width=2.5cm, minimum height=1cm, align=center}
    ]
        \node[entity, fill=blue!20] (browser) {Browser};
        \node[entity, fill=green!20, right=5cm of browser] (server) {Server};
        
        % Step 1
        \draw[-{Stealth}, thick] (browser) -- ++(0,1) -| node[above, pos=0.25] {1. GET /csrf} (server);
        
        % Step 2
        \draw[-{Stealth}, thick] (server) -- ++(0,-1) -| node[below, pos=0.25] {2. Set-Cookie: csrfToken=xxx} (browser);
        
        % Step 3
        \node[below=2cm of browser] (read) {Legge cookie (JavaScript)};
        
        % Step 4
        \draw[-{Stealth}, thick] (browser) -- ++(0,-3) -| node[above, pos=0.25] {4. POST + Header: X-CSRF-Token: xxx} (server);
        
        % Step 5
        \node[right=1cm of server, below=2cm] (validate) {Valida:\\Header == Cookie};
        
    \end{tikzpicture}
    \caption{Double-Submit Cookie Pattern}
\end{figure}

\section{Configurazione Cookie CSRF}

\begin{lstlisting}[caption={CSRF Cookie Options}]
{
    httpOnly: false,      // JavaScript deve leggerlo
    secure: true,         // HTTPS only (prod)
    sameSite: 'strict',   // No cross-origin
    maxAge: 24 * 60 * 60 * 1000  // 24 ore
}
\end{lstlisting}

%==============================================================================
% CAPITOLO 10: ENCRYPTION
%==============================================================================
\chapter{Crittografia Dati Sensibili}

\section{AES-256-GCM}

Dati sensibili (User-Agent, IP sessioni) sono criptati:

\begin{lstlisting}[caption={Encryption Service}]
const algorithm = 'aes-256-gcm';

function encryptString(plaintext) {
    const iv = crypto.randomBytes(16);
    const cipher = crypto.createCipheriv(
        algorithm, 
        ENCRYPTION_KEY, 
        iv
    );
    
    let encrypted = cipher.update(plaintext, 'utf8', 'hex');
    encrypted += cipher.final('hex');
    
    const authTag = cipher.getAuthTag();
    
    return iv.toString('hex') + ':' + authTag.toString('hex') + ':' + encrypted;
}
\end{lstlisting}

\section{Key Management}

\begin{warnbox}[Gestione Chiavi]
\begin{itemize}
    \item \texttt{DATA\_ENCRYPTION\_KEY}: 32 bytes, env variable
    \item Rotazione: Re-encrypt on access (lazy migration)
    \item Backup: HSM o KMS in produzione enterprise
\end{itemize}
\end{warnbox}

%==============================================================================
% CAPITOLO 11: ERROR HANDLING
%==============================================================================
\chapter{Gestione Errori}

\section{AppError Class}

Gerarchia errori personalizzati:

\begin{lstlisting}[caption={AppError Class Hierarchy}]
class AppError extends Error {
    constructor(message, statusCode, code, category, details) {
        super(message);
        this.statusCode = statusCode;
        this.code = code;           // UPPER_SNAKE_CASE
        this.category = category;    // AUTHENTICATION|AUTHORIZATION|...
        this.details = details;
        this.isOperational = true;   // true = expected error
        Error.captureStackTrace(this, this.constructor);
    }
    
    static unauthorized(message, code = 'UNAUTHORIZED') {
        return new AppError(message, 401, code, 'AUTHENTICATION');
    }
    
    static forbidden(message, code = 'FORBIDDEN') {
        return new AppError(message, 403, code, 'AUTHORIZATION');
    }
    
    static tooManyRequests(message, code = 'RATE_LIMIT') {
        return new AppError(message, 429, code, 'RATE_LIMIT');
    }
}
\end{lstlisting}

\section{Error Response Format}

\begin{lstlisting}[caption={Standard Error Response}]
{
    success: false,
    status: "error",
    error: {
        message: "Email o password non corretti",
        code: "UNAUTHORIZED",
        category: "AUTHENTICATION",
        suggestion: "Verifica le tue credenziali e riprova",
        details: {},
        requestId: "uuid-v4"
    },
    message: "Email o password non corretti"
}
\end{lstlisting}

%==============================================================================
% CAPITOLO 12: DEPLOYMENT CHECKLIST
%==============================================================================
\chapter{Deployment Security Checklist}

\section{Environment Variables}

\begin{longtable}{@{}lll@{}}
\toprule
\textbf{Variabile} & \textbf{Required} & \textbf{Note} \\
\midrule
\endhead
NODE\_ENV & Sì & production \\
JWT\_SECRET & Sì & Min 256 bit random \\
JWT\_ISSUER & Sì & silvi-api \\
JWT\_AUDIENCE & Sì & silvi-app \\
DATA\_ENCRYPTION\_KEY & Sì & 32 bytes hex \\
MONGODB\_URI & Sì & Replica set consigliato \\
REDIS\_URL & Opzionale & Per rate limiting distribuito \\
CORS\_ORIGIN & Sì & Domini whitelist \\
\bottomrule
\end{longtable}

\section{Security Headers}

Helmet.js configura automaticamente:

\begin{lstlisting}[caption={Helmet Configuration}]
{
    contentSecurityPolicy: true,
    crossOriginEmbedderPolicy: true,
    crossOriginOpenerPolicy: true,
    crossOriginResourcePolicy: true,
    dnsPrefetchControl: true,
    frameguard: { action: 'deny' },
    hidePoweredBy: true,
    hsts: { maxAge: 31536000, includeSubDomains: true, preload: true },
    ieNoOpen: true,
    noSniff: true,
    originAgentCluster: true,
    permittedCrossDomainPolicies: false,
    referrerPolicy: true,
    xssFilter: true
}
\end{lstlisting}

%==============================================================================
% APPENDICE
%==============================================================================
\appendix
\chapter{API Reference}

\section{Endpoints Authentication}

\begin{longtable}{@{}llll@{}}
\toprule
\textbf{Metodo} & \textbf{Endpoint} & \textbf{Auth} & \textbf{Rate Limit} \\
\midrule
\endhead
POST & /api/auth/register & No & General \\
POST & /api/auth/login & No & Auth (5/h) \\
POST & /api/auth/logout & Sì & General \\
POST & /api/auth/refresh & No & General \\
GET & /api/auth/check & Sì & General \\
GET & /api/auth/csrf & No & General \\
GET & /api/auth/profile & Sì & General \\
PATCH & /api/auth/profile & Sì & General \\
POST & /api/auth/password/change & Sì & Auth \\
POST & /api/auth/password/reset & No & Strict (3/h) \\
\midrule
POST & /api/2fa/setup & Sì & Strict \\
POST & /api/2fa/verify & Sì & Strict \\
POST & /api/2fa/disable & Sì & Strict \\
GET & /api/2fa/backup-codes & Sì & General \\
\bottomrule
\end{longtable}

\chapter{Database Schema}

\section{User Collection Indexes}

\begin{lstlisting}[caption={User Indexes}]
db.users.createIndex({ email: 1 }, { unique: true, sparse: true });
db.users.createIndex({ 'refreshTokens.jti': 1 });
db.users.createIndex({ 'trustedDevices.fingerprint': 1 });
db.users.createIndex({ createdAt: 1 });
\end{lstlisting}

\section{AuditLog Collection Indexes}

\begin{lstlisting}[caption={AuditLog Indexes}]
db.auditlogs.createIndex({ userId: 1, createdAt: -1 });
db.auditlogs.createIndex({ action: 1, createdAt: -1 });
db.auditlogs.createIndex({ ip: 1, createdAt: -1 });
db.auditlogs.createIndex({ createdAt: 1 }, { 
    expireAfterSeconds: 63072000  // 2 anni 
});
\end{lstlisting}

\end{document}
